\documentclass[10pt,a4paper]{amsart}
\usepackage[margin=19mm]{geometry}
\usepackage{amsmath,amssymb,mathtools}
\usepackage[scr=rsfs,cal=euler]{mathalfa}
\usepackage{xfrac}
\usepackage[unicode,hidelinks,bookmarks=false]{hyperref}
\newcommand{\ie}{i.e.\ }
\newcommand{\cf}{cf.\ }
\newcommand{\resp}{respectively\ }
\newcommand{\Subsection}{\subsection}
\providecommand{\nobreakdash}{\nobreak\hbox{-}}
\setlength{\emergencystretch}{2em}
\multlinegap0pt
\usepackage{amssymb}
\usepackage{bm}
\usepackage{mathtools}
\usepackage{tikz}
\usepackage[shortlabels]{enumitem}
\usepackage{float}
\newtheorem{lemma}{Lemma}
\title{Critical exponent for semilinear damped wave equations with weighted nonlinear terms and data from Sobolev spaces of negative order}
\author{Dinh Van Duong, Tuan Anh Dao}
\date{Source: arXiv:2508.07802v3 (CC BY 4.0)}
\begin{document}
\maketitle

\noindent\emph{Source and licence.} Critical exponent for semilinear damped wave equations with weighted nonlinear terms and data from Sobolev spaces of negative order. Version arXiv:2508.07802v3 (CC BY 4.0), distributed by arXiv under the Creative Commons Attribution 4.0 International licence, http://creativecommons.org/licenses/by/4.0/, as recorded in the arXiv licence metadata for this exact version and shown on the abstract page https://arxiv.org/abs/2508.07802v3. Full licence notice: \texttt{licenses/arxiv-2508.07802-CC-BY-4.0.txt}.\par\smallskip

\noindent\textit{Editorial scope.} Counted unit: the article's lemma LinearEstimates (source0.tex lines 425--431). The article's complete original proof of it is delivered with it (source0.tex lines 432--472, one \texttt{proof} environment of 41 lines). No mathematics is written, repaired or re-derived: the statement and the proof are the article's own lines, in the article's own notation, and no retained line is substituted or re-flowed.  No further source-local context is needed: every cross-reference of every retained line resolves inside the two delivered spans.  Nothing was omitted from any retained span, so the package carries no omission record.  No retained line contains a citation, so the wrapper carries no bibliography and no bibliography entry.  No retained line contains a figure environment, an image-inclusion command or drawing code, so no image is delivered and the package declares no asset dependency.  Editorial formatting only, in addition: an AMS \texttt{amsart} preamble that restates the article's own declarations and macros used by the retained lines, namely the environments \texttt{lemma} on separate counters, together with the standard packages those lines need.  The retained environments print in this document as Lemma 1 in order of appearance, whereas the article prints the same items with the numbering of its own full document.  The complete native source of the article as distributed in the arXiv e-print for this exact version is bundled unchanged as \texttt{sources/183052/source0.tex}.
\begin{lemma}\label{LinearEstimates}
    Let $n \geq 1, m \in (1,2], s \geq 0$ and $s+\beta \geq 0$.  Then, the following estimate holds for $j = 0,1$ and $t > 0$:
    \begin{align*}
        \| \partial_t^j \mathcal{K}(t,x)\ast_x \varphi(x)\|_{\dot{H}^s} 
        &\lesssim (1+t)^{-\frac{n}{2}(\frac{1}{m}-\frac{1}{2})-\frac{s+\beta}{2}-j} \|\varphi\|_{\dot{H}_m^{-\beta} \cap H^{s+j-1}}.
        \end{align*}
\end{lemma}

\begin{proof}
From the definition of the kernel $\mathcal{K}(t,x)$, we can get some pointwise estimates in the Fourier space, namely,
$$
|\partial_t^j \widehat{\mathcal{K}}(t,\xi)|\lesssim \begin{cases}
\mathrm{e}^{-ct}+|\xi|^{2j}\mathrm{e}^{-c|\xi|^2t}&\text{ for } |\xi|< \varepsilon^*,\\
\mathrm{e}^{-ct}&\text{ for } \varepsilon^* \leq|\xi|\leq N,\\
|\xi|^{j-1}\mathrm{e}^{-ct}&\text{ for } |\xi|> N,
\end{cases}
$$
with suitable constants $c>0,$  $\varepsilon^* \ll 1$ and $ N \gg 1$. Let $\chi_k= \chi_k(r)$ with $k\in\{\rm L,H\}$ be smooth cut-off functions having the following properties:
\begin{align*}
&\chi_{\rm L}(r)=
\begin{cases}
1 &\quad \text{ if }r\le \varepsilon^*/2 \\
0 &\quad \text{ if }r\ge \varepsilon^*
\end{cases}
\text{ and } \qquad
\chi_{\rm H}(r)= 1 -\chi_{\rm L}(r).
\end{align*}
Using Parseval's formula, we obtain
\begin{align*}
\|\partial_t^j \mathcal{K}(t,x)\ast_x \varphi(x)\|_{\dot{H}^s} =& \||\xi|^s \partial_t^j \widehat{\mathcal{K}}(t,\xi) \widehat{\varphi}(\xi)\|_{L^2} \\
\leq & \|\chi_{\rm L}(|\xi|)|\xi|^s \partial_t^j \widehat{\mathcal{K}}(t,\xi) \widehat{\varphi}(\xi)\|_{L^2} + \| \chi_{\rm H}(|\xi|)|\xi|^s \partial_t^j \widehat{\mathcal{K}}(t,\xi) \widehat{\varphi}(\xi)\|_{L^2}.  
\end{align*}
Concerning the small frequencies part,
applying H\"{o}lder's inequality with $1/m_0 + 1/m' = 1/2$ we gain
\begin{align*}
    \big\|\chi_{\rm L}(|\xi|)|\xi|^{s}\partial_t^j\widehat{\mathcal{K}}(t,\xi) \widehat{\varphi}(\xi)\big\|^2_{L^2} 
     &\leq \big\|\chi_{\rm L}(|\xi|)|\xi|^{s+\beta}\partial_t^j\widehat{\mathcal{K}}(t,\xi)\big\|_{L^{m_0}}^2 \||\xi|^{-\beta}\widehat{\varphi}(\xi)\|_{L^{m'}}^2.
\end{align*}
Thanks to the Hausdorff-Young's inequality, it holds $\||\xi|^{-\beta}\widehat{\varphi}(\xi)\|_{L^{m'}} \leq \|\varphi\|_{\dot{H}_m^{-\beta}}$ with $1/m' + 1/m = 1$  and $m \in (1,2]$. As a result, we obtain
\begin{align*}
     \big\|\chi_{\rm L}(|\xi|)|\xi|^{s}\partial_t^j\widehat{\mathcal{K}}(t,\xi) \widehat{\varphi}(\xi)\big\|_{L^2} &\lesssim \left(\big\|\chi_{\rm L}(|\xi|)|\xi|^{s+\beta+2j} e^{-c|\xi|^{2}(t+1)}\big\|_{L^{m_0}} + e^{-ct} \|\chi_{\rm L}(|\xi|) |\xi|^{s+\beta}\|_{L^{m_0}}\right) \|\varphi\|_{\dot{H}_m^{-\beta}}\\
     &\lesssim (1+t)^{-\frac{n}{2}(\frac{1}{m}-\frac{1}{2})-\frac{s+\beta}{2}-j} \|\varphi\|_{\dot{H}_m^{-\beta}}.
\end{align*}
Moreover, for the large frequencies part, we have immediately the following estimate:
$$
\|\chi_{\rm H}(|\xi|)|\xi|^{s}\partial_t^j \widehat{\mathcal{K}}(t,\xi)\widehat{\varphi}(\xi)\|_{L^2} \lesssim e^{-ct}\|\chi_{\rm H}(|\xi|)|\xi|^{s+j-1} \widehat{\varphi}(\xi)\|_{L^2} \lesssim e^{-ct} \|\varphi\|_{H^{s+j-1}}.
$$ 
The proof of Lemma \ref{LinearEstimates} has been completed.    
\end{proof}

\end{document}
